% !TEX root = ../main.tex
\chapter{Multiple Regression: Inference}
\label{ch:inference}

So far we have learned how to \emph{estimate} the parameters of a multiple regression model and what we can say about the estimators across repeated samples: under the Gauss--Markov assumptions OLS is unbiased and, among linear unbiased estimators, has the smallest variance. But estimation is only half of empirical work. Having computed a number like $\hb_1 = 0.083$ for the return to a year of schooling, the practitioner wants to know whether that number is reliably different from zero, whether it is compatible with a theoretical value like $0.10$, and how confident we should be in any of this. These are questions of \emph{inference}: statements about the unknown population parameters that come with an explicit probability of error.

To make such statements precise we need to know the full \emph{sampling distribution} of the OLS estimators, not just their first two moments. The Gauss--Markov assumptions pin down the mean and variance of $\hb_j$, but they say nothing about its shape; without a known distribution we cannot compute the probability that $\hb_j$ lands within a given distance of $\beta_j$. This chapter adds one final assumption --- that the error is normally distributed --- which delivers an exact normal sampling distribution and, through it, the $t$ and $F$ statistics that dominate applied econometrics. We develop the $t$ test for a single coefficient and its equivalence with confidence intervals, the trick that turns a test of a linear combination of parameters into an ordinary $t$ test, and the $F$ test for several restrictions at once. We close by separating two ideas that are constantly confused: \emph{statistical} significance, which is about precision, and \emph{economic} significance, which is about magnitude.

\section{The Normal Sampling Distribution}

Conditional on the sample values of the regressors, the sampling distribution of each $\hb_j$ is inherited entirely from the distribution of the unobserved errors. In Chapter~\ref{ch:mlr} we saw that $\hb_j$ is a linear function of the $y_i$, hence a linear function of the errors $u_i$; its mean and variance therefore followed from the mean and variance of $u$. To know its full distribution we must say something about the \emph{shape} of the error distribution. The simplest tractable choice, and the one that makes the finite-sample theory exact, is normality.

\asm{MLR.6 (Normality)}{
The population error $u$ is independent of the explanatory variables $x_1,x_2,\dots,x_k$ and is normally distributed with zero mean and variance $\sigma^2$:
\[
u \sim N(0,\sigma^2).
\]
}

This is a strong assumption. Independence of $u$ from the regressors is more than the zero conditional mean MLR.4 and the constant conditional variance MLR.5; it fixes the \emph{entire} conditional distribution of $u$, not just its first two moments. Consequently MLR.6 \emph{implies} both MLR.4 ($\E{u\given \vb x}=0$) and MLR.5 ($\Var{u\given \vb x}=\sigma^2$), and adds the assumption of a normal shape on top of them.

\defn{Classical Linear Model (CLM)}{
The assumptions MLR.1 through MLR.6 are called the \emph{classical linear model (CLM)} assumptions. Under the CLM, the conditional distribution of the dependent variable given the regressors is
\[
y \given \vb x \;\sim\; N\!\pr{\beta_0 + \beta_1 x_1 + \dots + \beta_k x_k,\ \sigma^2}.
\]
}

In words, the CLM says that for each fixed configuration of the regressors, $y$ is normally distributed about the population regression function $\beta_0+\beta_1 x_1+\dots+\beta_k x_k$ with a variance $\sigma^2$ that is the same at every configuration. The population regression function fixes the mean of this normal distribution; homoskedasticity fixes its spread; normality fixes its shape.

\rmkb[Why normality is plausible --- and when it fails]{The error $u$ is the sum of the many unobserved factors that influence $y$. If these factors are numerous, have comparable individual effects, and are roughly independent, a central limit argument suggests that their sum is approximately normal. The approximation is only as good as those premises: a single dominant factor, strongly skewed components, or strong dependence among them can all spoil it. Variables that are bounded, count-valued, or sharply skewed (wages, firm sizes) often have decidedly non-normal errors. This is exactly why Chapter~\ref{ch:asymptotics} develops large-sample inference, which dispenses with MLR.6 entirely: as $n\to\infty$ the OLS estimators are approximately normal regardless of the error distribution.}

Under the CLM we can strengthen what we know about OLS in two ways. First, its efficiency improves: under the Gauss--Markov assumptions OLS is best only within the class of \emph{linear} unbiased estimators (the BLUE property of Chapter~\ref{ch:mlr}); under the additional normality of MLR.6, OLS is the \emph{minimum variance unbiased estimator}, best among \emph{all} unbiased estimators, linear or not. Second, and central to this chapter, we obtain its exact sampling distribution.

\thm{Normal Sampling Distribution}{
Under the CLM assumptions MLR.1--MLR.6, conditional on the sample values of the independent variables,
\[
\hb_j \;\sim\; N\!\pr{\beta_j,\ \Var{\hb_j}}, \qquad j=0,1,\dots,k,
\]
and therefore the standardized estimator is standard normal:
\[
\frac{\hb_j - \beta_j}{\operatorname{sd}(\hb_j)} \;\sim\; N(0,1),
\qquad \operatorname{sd}(\hb_j)=\sqrt{\Var{\hb_j}}.
\]
}

\pf{
Each OLS estimator can be written as a linear combination of the errors,
\[
\hb_j = \beta_j + \sumi w_{ij}\, u_i,
\]
where the weights $w_{ij}$ are (nonrandom) functions of the sample values of the regressors alone --- for the slopes, $w_{ij}=\hat r_{ij}/\sum_{i} \hat r_{ij}^2$ in the partialling-out notation of Chapter~\ref{ch:mlr}. Conditional on the regressors the weights are constants. Under MLR.6 the $u_i$ are independent $N(0,\sigma^2)$ variables, and any linear combination of independent normal random variables is itself normal. Hence $\hb_j$ is normal; its mean $\beta_j$ and variance $\Var{\hb_j}$ are exactly the Gauss--Markov values derived in Chapter~\ref{ch:mlr}. Standardizing gives the $N(0,1)$ result.
}

The same argument shows more: \emph{any} linear combination of the $\hb_j$ is also exactly normal, because it too is a linear combination of the independent normal errors. We will exploit this repeatedly below.

\section{Testing a Single Parameter: The $t$ Test}

The standard normal result above is not yet usable, because $\operatorname{sd}(\hb_j)$ contains the unknown error variance $\sigma^2$. Replacing $\sigma$ by its estimator $\widehat\sigma$ (and hence $\operatorname{sd}(\hb_j)$ by the standard error $\operatorname{se}(\hb_j)$ from Chapter~\ref{ch:mlr}) changes the distribution from normal to Student's $t$, because we have introduced extra sampling noise through $\widehat\sigma$.

\thm{$t$ Distribution of the Standardized Estimator}{
Under the CLM assumptions MLR.1--MLR.6,
\[
\frac{\hb_j - \beta_j}{\operatorname{se}(\hb_j)} \;\sim\; t_{\,n-k-1},
\]
where $n-k-1$ is the residual degrees of freedom and $k+1$ is the number of parameters (intercept plus $k$ slopes) estimated in the model.
}

The degrees of freedom $n-k-1$ are exactly the divisor in the unbiased variance estimator $\widehat\sigma^2 = \tfrac{1}{n-k-1}\sumi \hat u_i^2$: we lose one degree of freedom for each of the $k+1$ estimated parameters. As $n-k-1$ grows the $t$ distribution approaches the standard normal, which is why large-sample $t$ tests are conducted using normal critical values.

\subsection{Carrying out the test}

To test a hypothesis about a single coefficient we form the $t$ statistic by plugging the hypothesized value into the standardized ratio. The most common null is that $\beta_j$ equals zero,
\[
H_0:\ \beta_j = 0,
\]
under which $x_j$ has no \emph{partial} effect on $y$ once the other regressors are controlled for. The corresponding statistic is simply the coefficient divided by its standard error,
\[
t_{\hb_j} = \frac{\hb_j}{\operatorname{se}(\hb_j)},
\]
which under $H_0$ follows $t_{\,n-k-1}$. We reject $H_0$ in favor of the alternative when $|t_{\hb_j}|$ (for a two-sided alternative) or $t_{\hb_j}$ (for a one-sided alternative) exceeds the relevant critical value, or equivalently when the $p$-value falls below the chosen significance level.

\rmkb[Reading the $t$ test correctly]{Several points trip up beginners and deserve to be stated once and for all.
\begin{itemize}
\item The null need not be zero. To test whether the coefficient equals some theoretical value $a_j$, use $H_0:\ \beta_j = a_j$ and the statistic $t = (\hb_j - a_j)/\operatorname{se}(\hb_j)$.
\item Match the tails. A two-sided alternative $H_1:\ \beta_j \ne a_j$ uses the absolute value of $t$ and a two-tailed critical value; a one-sided alternative ($H_1:\ \beta_j > a_j$ or $H_1:\ \beta_j < a_j$) uses the signed statistic and a one-tailed critical value. The $p$-value and critical value must be computed for the alternative actually being tested.
\item In large samples the $t$ test may be carried out using standard normal critical values.
\item ``Fail to reject $H_0$'' is not the same as ``accept $H_0$.'' Failing to reject means the data are consistent with $H_0$, not that $H_0$ is true; many other values of $\beta_j$ would also fail to be rejected.
\end{itemize}}

\subsection{Confidence intervals and the duality with testing}

A confidence interval inverts the test: instead of fixing one null value and asking whether the data reject it, it reports the entire set of null values the data would \emph{not} reject. Under the CLM the $(1-\alpha)$ confidence interval for $\beta_j$ is
\[
\hb_j \;\pm\; c_{\alpha/2}\cdot \operatorname{se}(\hb_j),
\]
where $c_{\alpha/2}$ is the two-tailed critical value from $t_{\,n-k-1}$ (for the usual $95\%$ interval and moderate degrees of freedom, $c_{\alpha/2}\approx 1.96$).

\state{Testing and confidence intervals are equivalent}{A value $a_j$ lies \emph{outside} the $(1-\alpha)$ confidence interval for $\beta_j$ if and only if the two-sided $t$ test of $H_0:\ \beta_j=a_j$ rejects at level $\alpha$. The interval is therefore a complete summary of every two-sided test at once: one can read off acceptance or rejection of any null directly, without recomputing a statistic.}

This duality is the reason confidence intervals are often preferred for reporting: a single interval communicates both the point estimate and the full range of values compatible with the data, and it answers every two-sided hypothesis test simultaneously.

\section{Testing a Linear Combination of Parameters}

Often the hypothesis of interest involves \emph{several} coefficients tied together by one linear relation --- for example, that the return to a year at a junior college equals the return to a year at a university, or that two elasticities are equal. The obstacle is that the relevant standard error is not printed by the regression software. A clever reparametrization sidesteps the difficulty by turning the question into an ordinary single-coefficient $t$ test.

\subsection{The reparametrization trick}

Consider the model
\[
y = \beta_0 + \beta_1 x_1 + \beta_2 x_2 + u,
\]
and suppose we wish to test a linear restriction relating the two slopes,
\[
H_0:\ \beta_1 = p + q\,\beta_2,
\]
for known constants $p$ and $q$. The idea is to express the restriction with a single new parameter $\theta$ that is zero exactly when $H_0$ holds. Write, with full generality,
\[
\beta_1 = p + q\,\beta_2 + \theta, \qquad\text{so that}\qquad H_0:\ \theta = 0.
\]
Substituting this into the regression equation and collecting terms,
\[
\ba
y &= \beta_0 + (p + q\,\beta_2 + \theta)\,x_1 + \beta_2 x_2 + u\\
&= \beta_0 + p\,x_1 + \theta\,x_1 + \beta_2(q\,x_1 + x_2) + u,
\ea
\]
which rearranges to
\[
\underbrace{y - p\,x_1}_{\text{new dependent variable}}
= \beta_0 + \theta\,x_1 + \beta_2 \underbrace{(q\,x_1 + x_2)}_{\text{new regressor}} + u.
\]
Now run OLS of the transformed dependent variable $y - p\,x_1$ on $x_1$ and the transformed regressor $q\,x_1 + x_2$. The coefficient on $x_1$ in this regression \emph{is} $\theta$, and the original null $H_0:\ \beta_1 = p+q\beta_2$ is now the elementary null $H_0:\ \theta=0$. We test it with the standard $t$ statistic $\hth/\operatorname{se}(\hth)$ that the software reports directly.

The payoff is twofold. The regression delivers $\hth = \hb_1 - p - q\,\hb_2$ (a point estimate of the linear combination) and, crucially, its standard error $\operatorname{se}(\hth) = \operatorname{se}(\hb_1 - q\,\hb_2)$ \emph{automatically} --- a quantity that would otherwise be awkward to compute by hand, because it requires the covariance between two estimated coefficients.

\ex[Equal returns to two-year and four-year college]{Let $\log(\text{wage}) = \beta_0 + \beta_1\,\text{jc} + \beta_2\,\text{univ} + \beta_3\,\text{exper} + u$, where $\text{jc}$ and $\text{univ}$ are years at a junior college and at a university. To test whether a year of each yields the same return, $H_0:\ \beta_1 = \beta_2$, take $p=0$ and $q=1$: define $\theta = \beta_1 - \beta_2$, regress $\log(\text{wage})$ on $\text{jc}$, the combined variable $\text{jc}+\text{univ}$, and $\text{exper}$, and read the $t$ statistic on $\text{jc}$. A negative, significant $\hth$ would say a junior-college year is worth less than a university year.}

\subsection{The covariance of two coefficients}

To see why the standard error of $\hb_1 - q\hb_2$ requires the covariance $\Cov{\hb_1}{\hb_2}$, expand the variance of the linear combination:
\[
\Var{\hb_1 - q\,\hb_2} = \Var{\hb_1} + q^2\,\Var{\hb_2} - 2q\,\Cov{\hb_1}{\hb_2}.
\]
The two individual variances are printed by the software, but the cross term is not. We compute it directly, using the fact (from Chapter~\ref{ch:mlr}) that each slope estimator is a linear function of the $y_i$ written through the partialling-out residuals $\hat r_{ij}$ --- the residuals from regressing $x_j$ on all the other regressors:
\[
\hb_1 = \frac{\sumi \hat r_{i1}\,y_i}{\sumi \hat r_{i1}^2},
\qquad
\hb_2 = \frac{\sumi \hat r_{i2}\,y_i}{\sumi \hat r_{i2}^2}.
\]
Because the observations are independently sampled, the $y_i$ are mutually independent (conditional on the regressors), each with variance $\Var{y_i\given \vb x} = \sigma^2$. Bilinearity of covariance then gives
\[
\ba
\Cov{\hb_1}{\hb_2}
&= \Cov{\frac{\sumi \hat r_{i1} y_i}{\sumi \hat r_{i1}^2}}{\frac{\sumi \hat r_{i2} y_i}{\sumi \hat r_{i2}^2}}\\[4pt]
&= \frac{1}{\pr{\sumi \hat r_{i1}^2}\pr{\sumi \hat r_{i2}^2}}\;
\Cov{\sumi \hat r_{i1} y_i}{\sumi \hat r_{i2} y_i}\\[4pt]
&= \frac{1}{\pr{\sumi \hat r_{i1}^2}\pr{\sumi \hat r_{i2}^2}}\;
\sumi \hat r_{i1}\,\hat r_{i2}\,\Var{y_i}\\[4pt]
&= \sigma^2\,\frac{\sumi \hat r_{i1}\,\hat r_{i2}}{\pr{\sumi \hat r_{i1}^2}\pr{\sumi \hat r_{i2}^2}},
\ea
\]
where the cross terms $\Cov{y_i}{y_j}=0$ for $i\ne j$ vanish by independence. The same idea gives a single tidy formula for any weighted combination of the two slopes. Writing $a_{ij} := \hat r_{ij}/\sum_l \hat r_{lj}^2$, so that $\hb_j = \sum_i a_{ij}\,y_i$, the combination $w_1\hb_1 + w_2\hb_2 = \sumi (w_1 a_{i1} + w_2 a_{i2})\,y_i$ is again a linear function of the independent $y_i$, hence
\[
\Var{w_1\hb_1 + w_2\hb_2}
= \sigma^2 \sumi \pr{w_1\,a_{i1} + w_2\,a_{i2}}^2.
\]
Expanding the square reproduces $w_1^2\Var{\hb_1} + w_2^2\Var{\hb_2} + 2w_1 w_2 \Cov{\hb_1}{\hb_2}$ with the variances and covariance just derived.
In practice one never plugs into this by hand; the reparametrization of the previous subsection lets OLS compute $\operatorname{se}(\hth)$ for us. The point of the derivation is conceptual: the standard error of a \emph{difference} of coefficients depends on how those coefficients covary, and ignoring the covariance term --- treating the two estimates as independent --- generally gives the wrong answer.

\rmkb[General linear combinations]{For a general linear combination $\sum_{j} w_j \hb_j$ of several coefficients, the variance is
\[
\Var{\textstyle\sum_j w_j \hb_j} = \sum_j w_j^2\,\Var{\hb_j} + \sum_{j\ne l} w_j w_l\, \Cov{\hb_j}{\hb_l},
\]
with each covariance computed exactly as above. Equivalently, in matrix form $\Var{\vb w' \hb} = \vb w'\,\Var{\hb}\,\vb w$, where $\Var{\hb}=\sigma^2 (\bX'\bX)^{-1}$ is the full covariance matrix of the coefficient vector. The reparametrization trick remains the practical route whenever the combination can be folded into a single redefined regressor.}

\section{Testing Multiple Linear Restrictions: The $F$ Test}

A single $t$ statistic tests one restriction. To test \emph{several} restrictions simultaneously --- ``do these three variables jointly belong in the model?'' --- we need a statistic that evaluates all the restrictions at once. This is the $F$ test.

\subsection{Restricted and unrestricted models}

A null that sets one population parameter to zero, say $H_0:\ \beta_k = 0$, is an \emph{exclusion restriction}: it excludes the corresponding variable. A null that constrains several parameters at once is a set of \emph{multiple restrictions}, and the test of such a null is a \emph{joint} (or multiple) hypothesis test. The leading case is a joint exclusion null,
\[
H_0:\ \beta_{k-q+1} = \beta_{k-q+2} = \dots = \beta_k = 0,
\]
which asserts that a designated group of $q$ regressors has no partial effect on $y$ as a group.

Imposing the restrictions gives the \emph{restricted model}, which drops the $q$ variables; the full model with all regressors is the \emph{unrestricted model}. The restricted model is \emph{nested} inside the unrestricted one --- it is the special case obtained by forcing $q$ coefficients to zero --- and the difference in the number of estimated parameters between the two models is exactly $q$, the number of restrictions being tested.

\subsection{The $F$ statistic in SSR form}

Dropping variables can never lower the sum of squared residuals: a restricted model fits no better than the model that is free to use the extra variables, so $\text{SSR}_r \ge \text{SSR}_{ur}$ always. The question is whether the \emph{increase} $\text{SSR}_r - \text{SSR}_{ur}$ is large enough --- relative to how well the unrestricted model fits --- to be incompatible with the restrictions being true. The $F$ statistic scales this increase appropriately.

\thm{The $F$ Statistic (SSR Form)}{
Under the CLM assumptions, with $q$ restrictions and $k$ regressors in the unrestricted model,
\[
F = \frac{(\text{SSR}_r - \text{SSR}_{ur})/q}{\text{SSR}_{ur}/(n-k-1)} \;\sim\; F_{\,q,\ n-k-1}
\qquad \text{under } H_0.
\]
}

The numerator is the increase in unexplained variation per restriction; the denominator is $\widehat\sigma^2$ from the unrestricted model, i.e.\ the unbiased estimate of the unrestricted error variance. The numerator degrees of freedom $q$ is the number of restrictions (equivalently, the difference in degrees of freedom between the two models), and the denominator degrees of freedom $n-k-1$ is the residual degrees of freedom of the unrestricted model. Because $\text{SSR}_r\ge\text{SSR}_{ur}$, the statistic is always nonnegative, and large values are evidence against $H_0$. We reject when $F$ exceeds the upper-tail critical value of the $F_{q,\,n-k-1}$ distribution.

\state{What rejection means}{If $F$ exceeds its critical value we reject $H_0$ and conclude that the excluded variables are \emph{jointly statistically significant} at the chosen level. If we cannot reject, the group is jointly insignificant --- the data do not establish that those variables, taken together, belong in the model. The SSR form is completely general: it is valid whether or not the restricted and unrestricted models share the same dependent variable, since it is built only from sums of squared residuals.}

\subsection{The $F$ statistic in $R^2$ form}

When the restricted and unrestricted models have the \emph{same} dependent variable --- as in any joint exclusion test --- the total sum of squares $\text{SST}$ is identical across the two models, so we may divide numerator and denominator by it and rewrite the statistic in terms of the $R^2$'s. Using $\text{SSR} = \text{SST}\,(1-R^2)$,
\[
F = \frac{(R_{ur}^2 - R_r^2)/q}{(1 - R_{ur}^2)/(n-k-1)} \;\sim\; F_{\,q,\ n-k-1}.
\]
This form is convenient because regression output reports $R^2$ directly. Note carefully that the denominator is $(1 - R_{ur}^2)$, the unexplained fraction of variation in the unrestricted model, divided by its degrees of freedom --- \emph{not} $(1-R_{ur})^2$. The $R^2$ form is restricted to the equal-dependent-variable case; whenever the dependent variable changes (e.g.\ a log transformation under one model), only the SSR form applies.

\subsection{Overall significance of the regression}

A special case worth singling out is the test that \emph{no} regressor matters,
\[
H_0:\ \beta_1 = \beta_2 = \dots = \beta_k = 0,
\]
the test of overall significance of the regression. Here the restricted model contains only the intercept, so $R_r^2 = 0$ and $q = k$. The $R^2$ form collapses to
\[
F = \frac{R^2/k}{(1 - R^2)/(n-k-1)} \;\sim\; F_{\,k,\ n-k-1}.
\]
This statistic appears in nearly every regression printout. Rejecting $H_0$ says only that \emph{at least one} regressor has explanatory power --- it does not identify which.

\subsection{$F$ versus $t$, and joint versus individual significance}

The $t$ and $F$ tests are intimately related but not interchangeable.

\state{Joint significance can survive when individual significance does not}{When two or more regressors are highly correlated (multicollinearity), each individual coefficient is estimated imprecisely, so the individual $t$ statistics may all be insignificant even though the variables clearly matter as a group. The $F$ test, which examines the variables \emph{jointly}, is essentially immune to this: it asks whether the group as a whole improves the fit, a question multicollinearity does not obscure. It is entirely possible --- and common --- for an $F$ test to reject the joint null while no single $t$ statistic is significant.}

The reason is geometric. Multicollinearity inflates the variance of each $\hb_j$ (recall $\Var{\hb_j} = \sigma^2/[\text{SST}_j(1-R_j^2)]$ from Chapter~\ref{ch:mlr}, which blows up as $R_j^2\to 1$) and induces a strong negative covariance between the correlated coefficients, so that each is poorly determined alone while their joint contribution is well determined. The $F$ statistic accounts for this covariance automatically; the separate $t$ statistics do not.

For a \emph{single} restriction the two tests coincide: the square of the $t$ statistic equals the $F$ statistic ($t^2 = F$ with $q=1$), and the two-sided $t$ test is exactly equivalent to the corresponding $F$ test. The $t$ statistic retains two practical advantages, however: it handles \emph{one-sided} alternatives, which the (inherently two-sided, nonnegative) $F$ statistic cannot, and it is more directly available, being printed automatically for every coefficient.

\subsection{General linear restrictions}

The $F$ machinery is not limited to exclusion restrictions. Any set of linear restrictions can be reduced to exclusion restrictions by the same reparametrization used for the single-restriction $t$ test. A general linear restriction has the form $\beta_1 = f(\boldsymbol\beta)$, where $f$ is a linear function of the other parameters. Writing it most generally as $\beta_1 = f(\boldsymbol\beta) + \theta$ and substituting into the equation, the restriction $H_0$ becomes $\theta = 0$:
\begin{itemize}
\item additive constants in $f$ get absorbed into a redefined dependent variable;
\item the linear coefficients in $f$ get absorbed into redefined regressors (combinations of the original ones).
\end{itemize}
After this transformation the joint null is a set of plain exclusion restrictions on the new parameters, and the ordinary SSR-form $F$ test applies. Thus exclusion restrictions are the canonical case to which all linear restrictions reduce.

\section{Economic versus Statistical Significance}

A coefficient can be statistically significant and economically trivial, or economically large and statistically uncertain. Keeping the two notions distinct is one of the most important habits in applied work.

\state{Two different questions}{The \emph{economic} (or \emph{practical}) significance of a variable is determined entirely by the \emph{size and sign} of its coefficient: how much does $y$ move, in units that matter, for a realistic change in $x_j$? The \emph{statistical} significance is determined entirely by the \emph{size of the $t$ statistic}: is the coefficient precisely enough estimated to rule out zero? A large coefficient with a small $t$ statistic is economically important but statistically shaky; a tiny coefficient with a huge $t$ statistic (common in very large samples) is statistically significant but may be economically negligible.}

Several practical guidelines follow.
\begin{itemize}
\item As the sample grows, smaller significance levels are warranted. Standard errors shrink with $n$, so even economically trivial effects eventually become statistically significant; using a stricter level offsets this and makes economic and statistical significance more likely to agree.
\item Check statistical significance first, but interpret economic significance in light of how the variable enters the equation --- its units, and whether it appears in levels, logs, or interactions, all change what ``the size of the coefficient'' means.
\item A variable with the wrong sign or an unexpected magnitude that is statistically insignificant can usually be set aside, but it may also be a symptom of a deeper specification problem --- an omitted variable, a wrong functional form, or multicollinearity --- and is worth investigating rather than ignoring.
\end{itemize}

\rmkb[Where this is heading]{We have added normality (MLR.6) to complete the classical linear model, obtained the exact normal sampling distribution of OLS, and built the $t$ and $F$ tests on top of it, along with the reparametrization device that handles any linear hypothesis. The exactness of these results rests on MLR.6, which is often only approximately true. Chapter~\ref{ch:asymptotics} shows that the same $t$ and $F$ procedures remain valid in large samples \emph{without} normality, by appeal to the central limit theorem; Chapter~\ref{ch:hetero} revisits inference once the homoskedasticity assumption MLR.5 is dropped.}
